\chapter{TÍCH HỢP CHATBOT KẾT HỢP TÌM KIẾM VÀ HỎI ĐÁP}\label{chap:chatbot}

Chatbot có hai chế độ do hệ thống tự chọn. \textbf{Chat} trò chuyện tự do, không truy xuất. \textbf{Research} chạy
Graph RAG và bắt buộc trích dẫn [n]. Khi người dùng vừa thêm PDF, bài đó trở thành \emph{bài đang hỏi} (focus):
các câu hỏi tiếp theo được trả lời từ chính bài đó, kèm các bài láng giềng trên đồ thị.

\section{Luồng một lượt trả lời}

\begin{enumerate}[label=(\arabic*)]
\item \textbf{Định tuyến:} đang có bài focus thì chọn Research; câu chào hỏi ngắn (biểu thức chính quy, tối đa
  8 từ) thì chọn Chat; còn lại chọn Research.
\item \textbf{Viết lại câu hỏi:} LLM biến câu hỏi cuối thành một truy vấn tiếng Anh độc lập, thay các đại từ
  (``bài này'', ``phương pháp đó'') bằng đối tượng cụ thể trong lịch sử hội thoại.
\item \textbf{Truy xuất Graph RAG} (Mục~\ref{sec:graphrag}): trả về danh sách nguồn đánh số, mỗi nguồn là một bài
  báo kèm các đoạn trích, vai trò (focus, tìm thấy, liên quan) và lý do liên quan.
\item \textbf{Sinh câu trả lời streaming} với prompt hệ thống yêu cầu: chỉ dùng ngữ cảnh, ghi [n] sau mỗi khẳng
  định, nói rõ khi không đủ thông tin, trả lời bằng ngôn ngữ của người hỏi.
\item \textbf{Lưu hội thoại} kèm nguồn và đồ thị con để mở lại được từ thanh bên.
\end{enumerate}

Mỗi lượt được phát thành chuỗi sự kiện SSE \code{meta} $\to$ \code{status} $\to$ \code{query} $\to$ \code{sources}
$\to$ \code{token}\ldots{} $\to$ \code{done}. Nhờ đó giao diện hiển thị tiến trình và nguồn trước khi câu trả lời
xuất hiện, rồi hiện câu trả lời dần từng chữ.

\section{Graph RAG}\label{sec:graphrag}

Graph RAG mở rộng RAG~\cite{rag} thông thường bằng một bước đi trên đồ thị tri thức, theo tinh thần của
GraphRAG~\cite{graphrag} và LightRAG~\cite{lightrag} nhưng ở mức đơn giản hơn.

\begin{itemize}
\item \textbf{Không có focus:} tìm kiếm lai lấy 5 bài ``hạt giống'' (mỗi bài 2 đoạn). Ba bài đầu được mở rộng sang
  2 láng giềng trên đồ thị, hoặc 4 nếu câu hỏi có ý ``liên quan / tương tự / so sánh / trích dẫn''.
\item \textbf{Có focus:} tóm tắt của bài cùng 8 đoạn khớp nhất trong chính bài (lọc theo \code{paper\_id} trên cả
  ChromaDB và BM25), cộng 4--8 bài liên quan của nó trên đồ thị.
\item \textbf{Quan hệ trên đồ thị} được viết thành câu (``A trích dẫn B'', ``A bàn về: \ldots'') và đưa vào ngữ
  cảnh. Nhờ vậy LLM trả lời được câu hỏi về quan hệ mà từng đoạn văn riêng lẻ không chứa.
\end{itemize}

Điểm liên quan của bài $q$ với bài $p$ được cộng dồn theo các đường đi 1--2 bước trên đồ thị:
\begin{equation}
\mathrm{score}(p,q) = \mathbb{1}[\mathrm{CITES}] + \cos_{\mathrm{SIMILAR}}(p,q)
  + \sum_{e \in N(p) \cap N(q)} \frac{w_{\mathrm{type}(e)}}{\log_2\bigl(2 + \deg(e)\bigr)},
\label{eq:lienquan}
\end{equation}
với $w = 0{,}6$ cho tác giả chung, $0{,}45$ cho khái niệm chung và $0{,}05$ cho chuyên mục chung. Thực thể càng phổ
biến (bậc $\deg(e)$ cao) thì càng ít giá trị phân biệt nên bị chia nhỏ điểm. Bài chỉ có chung chuyên mục thì không
được tính là liên quan. Mỗi kết quả kèm danh sách lý do để hiển thị, ví dụ ``trích dẫn bài này; cùng khái niệm
dense retrieval''.

\section{Mô hình ngôn ngữ}

\begin{table}[H]
\centering\small\setstretch{1.15}
\begin{tabularx}{\linewidth}{@{}P{0.2\linewidth}P{0.3\linewidth}L@{}}
\toprule
\textbf{Backend} & \textbf{Cấu hình} & \textbf{Ghi chú}\\
\midrule
vLLM tự host & Giao diện (dòng ``LLM'' ở thanh bên) hoặc \code{SL\_LLM\_PROVIDER=openai}, \code{SL\_LLM\_BASE\_URL} & Module \mo{llm\_server} host model Hugging Face, mặc định Qwen2.5-7B-Instruct-AWQ~\cite{qwen25,awq} trên 2 GPU T4\\
Claude & \code{SL\_LLM\_PROVIDER=anthropic} & Mặc định \code{claude-opus-5-5}; tự chuyển sang model dự phòng khi bị từ chối\\
API tương thích OpenAI khác & Gemini, Groq, Ollama & Dùng chung một client\\
Không LLM & \code{SL\_LLM\_PROVIDER=none} & Chế độ trích đoạn: chọn câu khớp nhất trong từng nguồn và liệt kê bài liên quan\\
\bottomrule
\end{tabularx}
\caption{Các backend LLM được hỗ trợ}\label{tab:llm}
\end{table}

vLLM~\cite{vllm} quản lý bộ nhớ đệm attention theo từng trang nhỏ (PagedAttention), ghép nhiều yêu cầu vào cùng
một lượt tính và chia mỗi lớp của mô hình cho 2 GPU cùng tính (tensor parallelism). Lượng tử hóa AWQ~\cite{awq}
nén trọng số xuống 4 bit, đưa mô hình 7B từ khoảng 15 GB xuống khoảng 5 GB để vừa với 2 GPU T4. Trên T4 (không hỗ
trợ bf16), module tự chọn \code{float16} và chế độ \code{enforce-eager} để tiết kiệm VRAM.

\section{Vận hành LLM từ xa}

LLM chạy trên máy GPU và được gọi qua một tunnel công khai, nên hệ thống phải chịu được việc tunnel đổi địa chỉ
hoặc mất kết nối: URL trycloudflare đổi mỗi lần khởi động lại tunnel, và phiên Kaggle có thể bị ngắt.

\begin{table}[H]
\centering\small\setstretch{1.15}
\begin{tabularx}{\linewidth}{@{}P{0.24\linewidth}L@{}}
\toprule
\textbf{Cơ chế} & \textbf{Cách làm}\\
\midrule
Đổi LLM khi đang chạy & Dòng ``LLM'' ở cuối thanh bên; bấm vào để nhập Base URL và API key. Ứng dụng kiểm tra kết nối qua \code{/v1/models}, tự thêm \code{/v1}, tự lấy tên model và áp dụng ngay, không cần khởi động lại. Cấu hình lưu ở \code{data/llm\_config.json}\\
Hiển thị trạng thái & Chấm xanh/đỏ và tên model trên dòng ``LLM'', làm mới mỗi 60 giây\\
Bảo vệ endpoint & vLLM bật \code{--api-key}; yêu cầu không có đúng khóa bị từ chối (HTTP 401). Khóa cố định đặt trong Kaggle Secrets\\
Tự quay về khi mất kết nối & Khi LLM không phản hồi, chatbot báo rõ lý do rồi trả lời bằng chế độ trích đoạn, nên không bao giờ trả về câu trả lời trống\\
Giữ phiên, tự mở lại tunnel & Trên máy GPU, \code{keep\_alive} kiểm tra vLLM và URL công khai mỗi 2 phút; sau 2 lần lỗi liên tiếp thì mở tunnel mới và chỉ in URL mới sau khi đã gọi thử thành công\\
\bottomrule
\end{tabularx}
\caption{Cơ chế vận hành LLM từ xa}\label{tab:vanhanh}
\end{table}

Khi kiểm thử với tunnel trycloudflare thật, URL mới cần khoảng 69 giây mới gọi được vì DNS chưa kịp cập nhật.
Kiểm thử này cũng phát hiện Cloudflare chặn User-Agent mặc định của thư viện \code{urllib}; nếu không gửi
User-Agent riêng, \code{keep\_alive} sẽ tưởng tunnel đã chết và đổi URL liên tục.
